\section{Longueur structurelle, action et réponse gravitationnelle}
\label{sec:vi-catalogo-dependencia-l-g}

Le catalogue réunit 324 trajets et 855 enregistrements métrologiques. Leur
interprétation suit la construction du même état HMT--MD, dans l'ordre causal
\[
 \mathrm{APP}\longrightarrow\mathrm{TRIT}\longrightarrow\mathrm{TPK}
 \longrightarrow X_{\mathrm{enr}}
 \longrightarrow\mathcal C^{\mathrm{disc}}_{\mathrm{joint}}
 \longrightarrow L_{\alpha}
 \longrightarrow(H\Lambda,R\Lambda)
 \longrightarrow G .
 \tag{A.1}\label{eq:vi-cat-cadena-lg}
\]
Les coordonnées \(\pi\), \(\varphi\), \(e\) et \(\alpha_{\mathrm{HMT}}\)
sont des sorties corrélées de l'image nonadique typée précédente ; elles ne
sont pas introduites comme valeurs conventionnelles pour sélectionner cette
réalisation. Pour \(N=12\cdot4\cdot120=5760\), la longueur marquée est
\[
 L_{\alpha}=\alpha_{\mathrm{HMT}}^{16}
 \frac{N-1}{4N}L_*
 =\frac{5759}{23040}\alpha_{\mathrm{HMT}}^{16}L_* .
 \tag{A.2}\label{eq:vi-cat-longitud-alpha}
\]
Dans cette section, on abrège \(\hbar=\hbar_{\mathrm{ret}}\) et
\(c=c_{\mathrm{int}}\). Sur la fibre positive non nulle, l'énergie \(H\),
la longueur d'action réduite \(\Lambda\) et la lecture radiale conjuguée
\(R\) satisfont les deux relations opératorielles suivantes. On définit
également \(M:=H/c^2\) et on déclare l'identification physique radiale
\(R:=GM/c^2\). Sa comparaison, sur la fibre inversible, avec la relation
opératorielle \(R=L_{\alpha}^{2}H/(\hbar c)\) donne
\[
 H\Lambda=\hbar c\,I,\qquad
 R\Lambda=L_{\alpha}^{2}I,\qquad
 R=\frac{L_{\alpha}^{2}}{\hbar c}H,\qquad
 \boxed{G=\frac{c^{3}L_{\alpha}^{2}}{\hbar}} .
 \tag{A.3}\label{eq:vi-cat-realizacion-lg}
\]
La reconnaissance métrologique de \(G\) est postérieure et n'intervient pas
dans le choix de l'état, du lecteur ou de ses coefficients.

\subsection{Code M10 et conservation de la loi de masses}
\label{sec:vi-catalogo-m10-activo}

Le code M10 du répertoire documentaire renvoie à la réalisation typée de
\eqref{eq:vi-cat-realizacion-lg}. M10 ne postule pas de particule gravitationnelle
et ne transforme pas en sortie HMT l'enregistrement externe conservé par
l'inventaire sous la classe \texttt{GRAVITATIONAL\_OR\_HYPOTHETICAL} : cet
enregistrement reste une donnée métrologique externe, séparée de la construction
\(L_{\alpha}\mapsto G\).

L'échelle énergétique induite par la même longueur est
\[
 Q_L:=\frac{\hbar c}{L_{\alpha}}.
 \tag{A.4}\label{eq:vi-cat-escala-q}
\]
Si des coordonnées équivalentes remplacent \(Q_L\) par \(Q_L'\), le coefficient
électronique se transforme conjointement selon
\[
 b_e'=b_e\frac{Q_L}{Q_L'},\qquad Q_L'b_e'=Q_Lb_e .
 \tag{A.5}\label{eq:vi-cat-cotransformacion-be}
\]
Les définitions de l'énergie de l'horloge et du facteur électronique donnent
\[
 E_{\mathrm{clk}}=Q_L,\qquad
 m_{\mathrm{clk}}=\frac{E_{\mathrm{clk}}}{c^2},\qquad
 b_e=\frac{E_e^{(0)}}{E_{\mathrm{clk}}},\qquad
 m_{\mathrm{clk}}b_e=\frac{E_e^{(0)}}{c^2}.
\]
L'annulation exacte de \(Q_L\) démontre que la réalisation de \(G\) ne dérive
pas de nouveau les masses, ne sélectionne pas les trajets et ne modifie pas
les tableaux déjà construits. La spécialisation \(\Xi_\Gamma=1\) du théorème
de compatibilité équivaut à \(\bar\lambda_\Gamma=R_{108}=L_{\alpha}\) : elle
identifie une coïncidence de longueurs et n'affirme ni ne crée l'existence
d'une espèce de masse de Planck dans le catalogue.

La décade \(10^{-34}\) appartient exclusivement à la branche absolue
action--horloge--masse, \(\hbar=\eta_{\mathrm{ret}}10^{-34}\mathcal U_S\) ;
elle n'est pas réinsérée dans l'exposant relatif, dans le caractère de trajet
ou comme correction par espèce.

\subsection{Canaux constitutifs et moment orienté de Catalan}
\label{sec:vi-catalogo-propietario-constitutivo}

Soient \(0<q_+,q_-<1\) les canaux étiquetés de transport et \(P_+,P_-\)
leurs projecteurs orthogonaux de feuillet. Pour \(s\in(0,1)\), on définit
\[
 \mathcal C_2(s)=\sum_{k\geq0}
 \frac{(-1)^k s^{2k+1}}{(2k+1)^2},\qquad
 \mathcal D=s\frac{\mathrm d}{\mathrm ds},\qquad
 f(s)=\frac{1-s^3}{1-s^4}.
 \tag{A.6}\label{eq:vi-cat-constitutivas}
\]
La convergence localement uniforme de la série et de ses dérivées permet
d'appliquer deux fois \(\mathcal D\). La série géométrique obtenue donne
\[
 \mathcal D^2\mathcal C_2(s)=\frac{s}{1+s^2},\qquad
 f(s)=\frac{1+s+s^2}{s(1+s)}\mathcal D^2\mathcal C_2(s).
\]
Les identités \(90=3\cdot30\) et \(120=4\cdot30\) déterminent
\[
 s_{\pm}=q_{\pm}^{30},\qquad
 r_{\pm}=f(s_{\pm})
 =\frac{1-q_{\pm}^{90}}{1-q_{\pm}^{120}}.
\]
Ainsi, \(V=r_+P_++r_-P_-\) réunit les réponses en conservant leurs étiquettes.
La série de \(\mathcal C_2\) converge absolument en \(s=1\) ; sa valeur est
\(G_{\mathrm{Cat}}\), distinguée du coefficient gravitationnel \(G\).

La fonction complète se retrouve à partir de \(f\) par
\[
 \mathcal C_2(s)=\int_0^s\log\!\left(\frac{s}{t}\right)
 \frac{1+t}{1+t+t^2}f(t)\,\mathrm dt.
\]
En effet, \((1+t)f(t)/(1+t+t^2)=1/(1+t^2)\) ; deux applications de
\(\mathcal D\) restituent \(s/(1+s^2)\). Les solutions de l'équation homogène
sont \(a\log s+b\) ; la condition \(\mathcal C_2(s)\to0\) lorsque
\(s\downarrow0\) fixe les deux coefficients à zéro. La réponse, sa restitution
régulière et le sens de la coordonnée de Catalan utilisée dans les fiches
sont ainsi explicités.
